\section{Primitive surplus derivations}
\label{app:surplus}

This appendix records the integration logic behind Proposition~\ref{prop:weak-cs}. Normalize \(\tau=1\) and suppress the common gross-utility term, which cancels from differences.

At the uniform source profile,
\[
a_u=1+s/3,
\qquad
b_u=1-s/3,
\qquad
x_u=1/2+s/6.
\]
If \(x_0\ge x_u\), realized uniform consumer surplus is
\begin{align*}
CS^u_{\rm sw}
={}&
\int_0^{x_u}(-a_u-z)\,dz
+
\int_{x_u}^{x_0}[-b_u-(1-z)-s]\,dz\\
&+
\int_{x_0}^{1}[-b_u-(1-z)]\,dz.
\end{align*}
If \(x_0<x_u\), the correct allocation has no switching:
\[
CS^u_{\rm ns}
=
\int_0^{x_0}(-a_u-z)\,dz
+
\int_{x_0}^{1}[-b_u-(1-z)]\,dz.
\]

For weak HBP, substitute the normalized prices in \eqref{eq:weak-hbp-prices}. The within-history cutoffs are
\[
x_A^d=\frac{2x_0+1+s}{6},
\qquad
x_B^d=\frac{2x_0+3-s}{6}.
\]
The primitive integral is
\begin{align*}
CS^d
={}&
\int_0^{x_A^d}(-p_A^d-z)\,dz
+
\int_{x_A^d}^{x_0}[-q_B^d-(1-z)-s]\,dz\\
&+
\int_{x_0}^{x_B^d}(-q_A^d-z-s)\,dz
+
\int_{x_B^d}^{1}[-p_B^d-(1-z)]\,dz,
\end{align*}
where all prices in this display are normalized net margins.

Subtracting the relevant uniform integral and restoring \(\tau\) gives Proposition~\ref{prop:weak-cs}.

\subsection{Sign on the switching branch}

With \(\tau=1\), the numerator of the switching-branch gap is
\[
N_1(x_0,s)
=
-52x_0^2+52x_0+1+36sx_0-34s+s^2.
\]
This is concave in \(x_0\). A switching subinterval inside weak dominance exists only for \(s\le3/5\). Its endpoint values satisfy
\[
N_1(x_u,s)
=
\frac{2(25s^2-72s+63)}9>0,
\]
and
\[
N_1(\bar x,s)
=
\frac{(1+s)(43-45s)}4>0.
\]
The first quadratic has negative discriminant and positive leading coefficient; the second is positive for \(s\le3/5\). Concavity therefore gives \(N_1>0\) throughout the switching part of the weak domain.

\subsection{Sign on the no-switch branch}

The no-switch numerator is
\[
N_0(x_0,s)
=
s^2+12sx_0-14s-8x_0^2+8x_0+5,
\]
also concave in \(x_0\). At the lower endpoint,
\[
N_0(1/2,s)=(1-s)(7-s)>0.
\]
If \(s\le3/5\), the relevant upper endpoint is \(x_u\), where
\[
N_0(x_u,s)
=
\frac{25s^2-72s+63}{9}>0.
\]
If \(s>3/5\), the weak upper endpoint is \(\bar x<x_u\), where
\[
N_0(\bar x,s)
=
\frac{(1-s)(5s+13)}2>0.
\]
Hence the weak profile consumer-surplus gap is positive throughout the source weak-dominance domain.

\subsection{Exact regression points}

Two exact rational checks are retained. At \((s,x_0)=(99/100,501/1000)\), the actual uniform allocation is no-switch and
\[
CS^d-CS^u=\frac{17993}{4500000}>0.
\]
At \((s,x_0)=(1/10,7/10)\), which lies on the one-way-switch branch and inside the corrected pure-overlap domain,
\[
CS^d-CS^u=\frac{221}{720},
\]
while the accepted Eq.~(15) expression evaluates to \(1279/3600\). These checks distinguish the branch-ordering problem from the independent algebraic discrepancy in the accepted display.
